I cannot honestly give a complete proof of the full statement from only the facts explicitly listed in the prompt.

What I **can** give rigorously is a substantial partial solution:

1. the standard closure and parity relations for the translate walk;
2. a complete proof that the triangle case \(n=3\) is impossible;
3. the easy reduction that if there are only two edge-direction classes, then \(P\) is a parallelogram.

The remaining gap is the global \(n\ge 4\) step: from “at least three direction classes” one still needs a rigorous local-cell obstruction, and that uses more structure than is explicitly stated in the prompt.

---

Let
\[
x_i:=s_i+t_i \qquad (1\le i\le n).
\]

By the exposed-point property, each \(x_i\) is a uniquely represented sum, so
\[
\sigma(x_i)=\sigma(t_i)\neq 0.
\]
Hence there are at least \(n\) nonzero points. Since by hypothesis
\[
|\{p:\sigma(p)\neq 0\}|=n,
\]
it follows that the nonzero points are exactly
\[
\{x_1,\dots,x_n\}.
\]

So any point \(p\notin \{x_1,\dots,x_n\}\) must satisfy
\[
\sigma(p)=0.
\]

## 1. Closure and parity

From
\[
t_{i+1}=t_i+m_i a_i \qquad (i \bmod n),
\]
summing around the cycle gives
\[
\sum_{i=1}^n m_i a_i = 0.
\]

Also, because the signs alternate along
\[
F_i=\{t_i,t_i+a_i,\dots,t_i+m_i a_i\},
\]
we have
\[
\sigma(t_{i+1}) = (-1)^{m_i}\sigma(t_i).
\]
Multiplying these relations around the cycle,
\[
\sigma(t_1)=(-1)^{m_1+\cdots+m_n}\sigma(t_1),
\]
so
\[
m_1+\cdots+m_n \equiv 0 \pmod 2.
\]

Thus we have proved:

\[
\boxed{\sum_{i=1}^n m_i a_i=0,\qquad \sum_{i=1}^n m_i\equiv 0\pmod 2.}
\]

## 2. Triangle case

For the triangle case, I will use the standard stronger conclusion from the earlier subproblems that every translate lies on one of the three edge chains:
\[
T=F_1\cup F_2\cup F_3.
\]
(Your hint also uses this.)

Assume now that \(n=3\). Then
\[
a_1+a_2+a_3=0,
\]
and \(a_1,a_2\) are linearly independent.

From closure,
\[
m_1a_1+m_2a_2+m_3a_3=0.
\]
Substitute \(a_3=-a_1-a_2\):
\[
(m_1-m_3)a_1+(m_2-m_3)a_2=0.
\]
By linear independence,
\[
m_1=m_2=m_3=:m.
\]

From parity,
\[
3m\equiv 0 \pmod 2,
\]
hence \(m\) is even.

Also \(m\neq 0\), because \(m=0\) would give
\[
t_1=t_2=t_3,
\]
so \(T\) would be a singleton, impossible since both signs occur. Therefore
\[
m\ge 2.
\]

### Case 1: \(m=2\)

Take
\[
p:=s_3+t_1+a_1.
\]
Then
\[
p-s_3=t_1+a_1\in F_1,
\]
\[
p-s_2=t_1+a_1+a_2=t_3+a_3\in F_3,
\]
and
\[
p-s_1=t_1+2a_1+a_2=t_2+a_2\in F_2.
\]
So all three vertices contribute.

Their signs are:
\[
\sigma(t_1+a_1)=-\sigma(t_1),
\]
\[
\sigma(t_3+a_3)=-\sigma(t_3)=-\sigma(t_1),
\]
\[
\sigma(t_2+a_2)=-\sigma(t_2)=-\sigma(t_1),
\]
because \(m=2\) is even, hence
\[
\sigma(t_2)=\sigma(t_3)=\sigma(t_1).
\]
Therefore
\[
\sigma(p)=-3\,\sigma(t_1)\neq 0.
\]

We must check that \(p\) is not one of the three lonely points.

- If \(p=x_1=s_1+t_1\), then
  \[
  s_3-s_1+a_1=0,
  \]
  i.e.
  \[
  (a_1+a_2)+a_1=0,
  \]
  so
  \[
  2a_1+a_2=0,
  \]
  impossible.

- If \(p=x_2=s_2+t_2\), then
  \[
  s_3+t_1+a_1=s_2+t_1+2a_1,
  \]
  so
  \[
  a_1+a_2+a_1=a_1+2a_1,
  \]
  hence
  \[
  a_2=a_1,
  \]
  impossible.

- If \(p=x_3=s_3+t_3\), then
  \[
  t_1+a_1=t_3=t_1+2a_1+2a_2,
  \]
  so
  \[
  a_1+2a_2=0,
  \]
  impossible.

So \(p\notin \{x_1,x_2,x_3\}\), yet \(\sigma(p)\neq 0\), contradiction.

### Case 2: \(m\ge 4\)

Take
\[
p:=s_3+t_1+2a_1.
\]
Then
\[
p-s_3=t_1+2a_1\in F_1.
\]

Now check the other two possible translates.

For \(s_2\):
\[
p-s_2=t_1+2a_1+a_2.
\]
This is not in \(F_1\), because that would force \(a_2\) to be a multiple of \(a_1\).
It is not in \(F_2\), because
\[
t_1+2a_1+a_2=t_2+ja_2=t_1+ma_1+ja_2
\]
would imply
\[
(2-m)a_1+(1-j)a_2=0,
\]
so \(2=m\), impossible.
It is not in \(F_3\), because
\[
t_1+2a_1+a_2=t_3+ja_3=t_1+ma_1+ma_2-j a_1-j a_2
\]
would imply
\[
(2-m+j)a_1+(1-m+j)a_2=0,
\]
whose two coefficients differ by \(1\), impossible.

Hence
\[
p-s_2\notin T.
\]

For \(s_1\):
\[
p-s_1=t_1+3a_1+a_2.
\]
Again this is not in \(F_1\).
It is not in \(F_2\), because
\[
t_1+3a_1+a_2=t_1+ma_1+ja_2
\]
would force \(3=m\), impossible.
It is not in \(F_3\), because
\[
(3-m+j)a_1+(1-m+j)a_2=0
\]
has coefficients differing by \(2\), impossible.

So
\[
p-s_1\notin T.
\]

Therefore the only contribution to \(p\) is from \(s_3\), namely
\[
\sigma(p)=\sigma(t_1+2a_1)=\sigma(t_1)\neq 0.
\]

Again \(p\notin \{x_1,x_2,x_3\}\):

- \(p=x_1\) would give
  \[
  s_3-s_1+2a_1=0 \implies 3a_1+a_2=0,
  \]
  impossible.

- \(p=x_2=s_2+t_2\) would give
  \[
  a_1+a_2+2a_1=a_1+ma_1,
  \]
  hence
  \[
  a_2=(m-2)a_1,
  \]
  impossible.

- \(p=x_3=s_3+t_3\) would give
  \[
  2a_1=ma_1+ma_2,
  \]
  impossible.

So again we get a contradiction.

Hence:

\[
\boxed{n=3\text{ is impossible}.}
\]

## 3. If there are only two direction classes, then \(P\) is a parallelogram

A convex polygon cannot have two consecutive parallel edges. So if there are only two edge-direction classes, the classes must alternate around the boundary. A given direction class can occur at most twice in a convex polygon (once in each orientation around the cycle). Therefore there are exactly four edges, and opposite edges are parallel. Hence \(P\) is a parallelogram.

So:

\[
\boxed{\text{If }P\text{ has only two edge-direction classes, then }n=4\text{ and }P\text{ is a parallelogram}.}
\]

## Conclusion

From the explicitly stated material, I can rigorously conclude:

- the translate walk satisfies
  \[
  \sum_i m_i a_i=0,\qquad \sum_i m_i\equiv 0\pmod 2;
  \]
- triangles are impossible;
- if there are only two direction classes, then the only possible polygon is a parallelogram.

So the full theorem is reduced to the remaining case
\[
n\ge 4,\qquad P\text{ has at least three direction classes}.
\]

I do **not** see a completely rigorous derivation of the final contradiction for that remaining case from only the two listed prior facts. The missing ingredient is a further local-cell obstruction controlling how the chains \(F_i\) can interact when three direction classes are present.

If you want, I can next do one of two things:

1. try to reconstruct that missing local-cell lemma and push the proof further; or  
2. write the above as a polished “partial solution” in theorem-proof style suitable for submission.